\documentclass[11pt]{letter}
\usepackage[margin=1in]{geometry}
\signature{Albert Linda\\Department of Materials Science and Engineering\\
Indian Institute of Technology Kanpur, Kanpur 208016, India\\
albert.hzbn@gmail.com}
\address{}
\begin{document}
\begin{letter}{Editor\\Journal of Applied Physics}
\opening{Dear Editor,}

I submit the manuscript ``Special optimal thermal cells: D-optimal thermal
displacements for finite-temperature effective-harmonic phonons'' for
consideration as a regular article in the Journal of Applied Physics.

Finite-temperature phonons underlie applied predictions of phase stability,
vibrational thermodynamics, and thermal transport, and the number of
first-principles force calculations is the main cost of obtaining them. The
manuscript asks whether designed, rather than randomly sampled, supercell
displacements can reduce that cost, and answers it with fully matched
controls. For eleven crystals of five structure classes (fcc Al, Ag, Au, Cu;
bcc Mo, W; rocksalt NaCl, MgO, PbTe; diamond Si; hcp Mg) the designs are
compared with converged large-cell references, with effective force
constants fitted to canonical ensembles in the same cell, and, at identical
numbers of force calls, cells, electronic settings, and fits, with the
established alternatives: canonical sampling with TDEP-type least-squares
fitting, its iterated form sTDEP, TDEP fitted to \textit{ab initio} molecular
dynamics (Al), SSCHA (with DFT forces for Al and on DFT-fitted anharmonic model
potentials), special displacements, and thermal rattle with hiPhive regression.

The principal results are:
\begin{itemize}
\item a D-optimal antithetic design (SOTC-D) with an estimator theory: the
effective force constants are the Gaussian-averaged Hessian, SOTC-D estimates
them with an $O(1/M)$ bias at fourth order, and inversion pairs cancel cubic
forces exactly;
\item with the same trial, SOTC-D has the lowest dispersion error of five
designs in 10 of 11 crystals at two and at eight force calls; its
bootstrapped form beats one-shot canonical sampling in all
eleven crystals at every budget from 4 to 16 calls (supercell basis) and two-iteration sTDEP
with DFT forces in all eleven at eight calls, and on anharmonic
model potentials of five crystals its fully self-consistent form reaches a
given accuracy with a factor of 2--16 fewer force calls than SSCHA in 13 of
14 cases;
\item regression method (nine least-squares, sparse, and Bayesian fits)
matters little, whereas a crystal-symmetry basis reduces the parameter count
of the cubic crystals by up to a factor of 30 at no loss of accuracy;
\item a quantitative separation of sampling error from finite-temperature
renormalization (up to 1.4~THz) and supercell truncation (up to 0.5~THz),
room-temperature Debye--Waller factors within 6\% of tabulated values
(converged references for Al, Cu, and Au; SOTC-D models within 2--11\%),
and non-analytic LO--TO corrections for the polar crystals;
\item on a volume grid, a reproducible thermal expansion of Al within 3\% of
experiment at 300~K and, with a cubic free-energy fit, 7--8\% below it at 400~K, with a statement of where the small budgets of the
design stop being sufficient (the free-energy curvature at higher
temperatures).
\end{itemize}
All inputs, outputs, and analysis scripts are openly available.

The manuscript is not under consideration elsewhere. Suitable reviewers
include experts in first-principles anharmonic lattice dynamics and
force-constant fitting methods.

\closing{Sincerely,}
\end{letter}
\end{document}
